\documentclass[aps,prb,reprint,amsmath,amssymb,superscriptaddress,nofootinbib]{revtex4-2}

% --- PACKAGES ---
\usepackage{graphicx}
\usepackage{array}
\usepackage{dcolumn}
\usepackage{bm}
\usepackage{physics}
\usepackage{hyperref}
\usepackage{microtype}
\usepackage{xcolor}
\usepackage{booktabs}

\hypersetup{
    colorlinks=true,
    linkcolor=blue,
    citecolor=blue,
    urlcolor=blue
}

% --- MODEL LABELS (change here to rename everywhere) ---
\newcommand{\mS}[1]{\ensuremath{\textsf{S#1}}}
\newcommand{\mQ}{\ensuremath{\textsf{Q}}}
\newcommand{\mQb}{\ensuremath{\textsf{Q}'}}
\newcommand{\mL}{\ensuremath{\textsf{L}}}
\newcommand{\mU}{\ensuremath{\textsf{U}}}
\newcommand{\mQJ}{\ensuremath{\textsf{Q+J}_1}}
\newcommand{\mUJ}{\ensuremath{\textsf{U+J}_1}}
\newcommand{\mFJ}{\ensuremath{\textsf{F+J}_1}}
\newcommand{\mD}{\ensuremath{\textsf{D}}}
\newcommand{\mDU}{\ensuremath{\textsf{D}_{\textsf{U}}}}
\newcommand{\mDJ}{\ensuremath{\textsf{D+J}_1}}
% Figures live in paper/figures/. A missing file gives a framed placeholder
% instead of an error, so the manuscript always compiles.
\newcommand{\figfile}[2]{%
  \IfFileExists{figures/#2}{\includegraphics[width=#1]{figures/#2}}{%
    \fbox{\parbox{0.95#1}{\centering\small\vspace{2em}%
      Missing figure: \texttt{\detokenize{#2}}\vspace{2em}}}}}
\newcommand{\Eh}{\,E_h}
\newcommand{\mHa}{\,\mathrm{mHa}}

\begin{document}

% --- TITLE & METADATA ---
\title{What a Neural Jastrow Factor Can and Cannot Learn:\\
A Controlled Variational Monte Carlo Study of the Beryllium Atom}

\author{Francesco Agus}
% Matricola rimossa: il repository e' pubblico. Inserisci qui
% l'affiliazione istituzionale, oppure lascia la riga commentata.
%\affiliation{}

\date{\today}

% --- ABSTRACT ---
\begin{abstract}
We study the ground state of the beryllium atom with variational Monte Carlo
(VMC) and fixed-node diffusion Monte Carlo (FN-DMC). The trial wave functions
are products of a determinant part and a Jastrow factor that contains a
neural network. The Jastrow factor is positive, so FN-DMC measures the quality
of the determinant part alone. In a series of four single-determinant trial
functions with the same node, a neural network must learn the
electron-electron cusps. It reproduces the antiparallel-spin cusp to $0.7\%$
and misses the parallel-spin cusp by $81\%$, with no measurable effect on the
energy. For a
complete-active-space expansion over the $2s^2$ and $2p^2$ configurations we
show that the determinant part factorises into a positive function and a
nodal function of two radial functions and one number, and we give the
transformations of the orbitals that leave the node unchanged. Across nine
trial functions, those with a radial modulation of the $1s$ orbital have a
fixed-node energy $0.20(3)\mHa$ lower than those without it. Neither the VMC
energy nor the configuration-mixing coefficient detects this difference, and
a single gauge-invariant number does not order nodes whose radial functions
differ. A $3d^2$ configuration and a one-body Jastrow term do not change the
node measurably. The difference $E_{\mathrm{VMC}} - E_{\mathrm{FN}}$ is
$0.2$--$0.8\mHa$, larger than the fixed-node error, which is $0.14(2)\mHa$ for
the trial functions with a modulated $1s$ orbital.
\end{abstract}

\maketitle


% ============================================================================
\section{Introduction}
\label{sec:intro}
% ============================================================================

A common trial wave function for electrons has the Slater-Jastrow form,
\begin{equation}
    \Psi(\mathbf{R}) = D(\mathbf{R})\, e^{J(\mathbf{R})} .
    \label{eq:slater_jastrow}
\end{equation}
The determinant part $D$ is antisymmetric. The Jastrow factor $e^{J}$ is
symmetric and positive, and describes the correlation of electron pairs at
short distance. Since $e^{J} > 0$, the nodal surface of $\Psi$, the set where
$\Psi = 0$, is the nodal surface of $D$. Neural networks are now used to
represent $J$ and, in some approaches, the whole wave
function~\cite{carleo2017,pfau2020,hermann2020}. A network inside $J$
increases the flexibility of the positive factor. It does not change the node.

The variational Monte Carlo (VMC) energy contains two errors: the error of the
node and the error of the positive factor. Fixed-node diffusion Monte Carlo
(FN-DMC) separates them. In the limit of zero time step it gives the lowest
energy compatible with the node of the trial function, whatever the positive
factor~\cite{reynolds1982,foulkes2001}. The difference
$E_{\mathrm{VMC}} - E_{\mathrm{FN}}$ is then the error of the positive factor,
and $E_{\mathrm{FN}} - E_{\mathrm{exact}}$ is the error of the node. When two
trial functions differ in both factors, their VMC energies do not tell which
factor is responsible for the difference.

The beryllium atom is a test case for this question. Its $2s^2$ configuration
is nearly degenerate with $2p^2$, and a single determinant gives a fixed-node
error of order $10\mHa$~\cite{todobehfnode}. The exact non-relativistic
energy is known to better than $1\,\mu E_h$~\cite{pachucki2006}. With four
electrons the node of a few-configuration wave function can be written in
closed form~\cite{bressanini_nodes}.

We study two sets of trial functions. The first is a series of four
single-determinant trial functions with the same node and different Jastrow
factors, from none to a neural network that must learn the electron-electron
cusps. The second is a complete-active-space expansion over $2s^2$ and $2p^2$,
with optional $3d^2$ configurations, multiplied by a Jastrow factor with a
three-body neural network. For the second set we show that the determinant
part factorises into a positive function and a nodal function, and we list
the transformations of the orbitals that leave the node unchanged. We then
compute VMC and FN-DMC energies for nine trial functions that differ in one
element at a time: the radial form of the orbitals, a one-body Jastrow term,
the $3d^2$ configurations, and the random seed of the optimisation.

The results are the following. In the single-determinant series the network
learns the antiparallel-spin cusp and not the parallel-spin cusp, and the
energy does not depend on the latter. The VMC error of the multi-configuration
trial functions comes mostly from the Jastrow factor. The node depends on the
shape of the $1s$ orbital: the trial functions in which the $1s$ orbital is
modulated have a fixed-node energy $0.20(3)\mHa$ lower than the others.
Neither the VMC energy nor the mixing coefficient of the $2p^2$ configurations
shows this. The mixing coefficient is not identifiable from the optimisation,
and its gauge-invariant version does not order nodes whose radial functions
differ. The $3d^2$ configurations and a one-body Jastrow term do not change
the node measurably.

Section~\ref{sec:wavefunction} defines the trial functions.
Section~\ref{sec:nodes} derives the structure of the nodal surface.
Section~\ref{sec:methods} describes the VMC and FN-DMC calculations.
Section~\ref{sec:results} presents the results, Sec.~\ref{sec:discussion}
discusses them, and Sec.~\ref{sec:limitations} states the limitations of the
study.


% ============================================================================
\section{Trial wave functions}
\label{sec:wavefunction}
% ============================================================================

\subsection{Hamiltonian and antisymmetry}

In Hartree atomic units the Hamiltonian of the beryllium atom with a fixed
nucleus of charge $Z = 4$ is
\begin{equation}
    \hat{H} = -\frac{1}{2}\sum_{i=0}^{3} \nabla_i^2
              - \sum_{i=0}^{3} \frac{Z}{r_i}
              + \sum_{0\le i<j\le 3} \frac{1}{r_{ij}} ,
    \label{eq:hamiltonian}
\end{equation}
with $r_i = |\mathbf{r}_i|$ and $r_{ij} = |\mathbf{r}_i - \mathbf{r}_j|$.
$\hat H$ does not act on spin. We fix $S_z = 0$, with electrons $0, 1$ of
spin up and $2, 3$ of spin down. For two spatial orbitals $\phi_a$, $\phi_b$
the determinant part is the product of two $2\times2$ determinants,
\begin{equation}
    D[\phi_a, \phi_b] = \det M_\uparrow \, \det M_\downarrow ,
    \label{eq:det_product}
\end{equation}
with $M_\uparrow = \bigl(\phi_k(\mathbf{r}_i)\bigr)$ for $i \in \{0, 1\}$ and
$M_\downarrow = \bigl(\phi_k(\mathbf{r}_i)\bigr)$ for $i \in \{2, 3\}$,
$k \in \{a, b\}$.
It changes sign under the exchange of two electrons of the same spin. Of the
six electron pairs, two have parallel spins, $(0,1)$ and $(2,3)$, and four
antiparallel spins.

\subsection{Orbitals}
\label{ssec:orbitals}

When an electron approaches the nucleus, $-Z/r_i$ diverges. A finite local
energy requires, for an $s$ orbital, the electron-nucleus cusp
condition~\cite{kato1957}
\begin{equation}
    \left. \frac{\phi'(r)}{\phi(r)} \right|_{r\to 0} = -Z .
    \label{eq:en_cusp}
\end{equation}
We impose it exactly. For $\phi_{1s} \propto e^{-\alpha_{1s} r}$ it gives
$\alpha_{1s} = Z$, which we keep fixed. For
$\phi_{2s} \propto (r - c)\, e^{-\alpha_{2s} r}$ one has $\phi(0) = -c$ and
$\phi'(0) = 1 + \alpha_{2s} c$, so Eq.~\eqref{eq:en_cusp} fixes the position of
the radial node,
\begin{equation}
    c = \frac{1}{Z - \alpha_{2s}} .
    \label{eq:c_from_cusp}
\end{equation}
$c$ is recomputed from $\alpha_{2s}$ at every evaluation. The $2p$ and $3d$
orbitals vanish at the nucleus and carry no cusp condition.

Each orbital is multiplied by a radial modulation built on the coordinate
\begin{equation}
    t(r) = \frac{r^2}{1+r} , \qquad t(0) = 0 , \quad t'(0) = 0 .
    \label{eq:t_coordinate}
\end{equation}
The orbitals are
\begin{align}
    \phi_{1s}(r)  &= e^{-Z r}\, m_{1s}(r) ,
    \label{eq:phi1s} \\
    \phi_{2s}(r)  &= (r - c)\, e^{-\alpha_{2s} r}\, m_{2s}(r) ,
    \label{eq:phi2s} \\
    \phi_{2p_k}(\mathbf{r}) &= x_k\, e^{-\alpha_{2p} r}\, m_{2p}(r) ,
    \qquad k \in \{x, y, z\} ,
    \label{eq:phi2p}
\end{align}
and, for the $3d$ shell, the five components of
$T(\mathbf{r})\, e^{-\alpha_{3d} r}\, m_{3d}(r)$, with
$T_{ab}(\mathbf{r}) = x_a x_b - \delta_{ab} r^2/3$. The modulation is
\begin{equation}
    m_{nl}(r) = 1 + c^{(1)}_{nl}\, t(r) + c^{(2)}_{nl}\, t(r)^2 .
    \label{eq:modulation}
\end{equation}
Since $t(0) = t'(0) = 0$, $m_{nl}$ has unit value and zero slope at the
nucleus, and Eq.~\eqref{eq:en_cusp} holds for any value of the coefficients.
The coefficients are bounded, $c^{(1)} \in [-0.15, 2.5]$ and
$c^{(2)} \in [-0.15, 0.8]$, and the exponents are bounded,
$\alpha_{2s}, \alpha_{2p} \in [0.3, 3]$ and $\alpha_{3d} \in [0.3, 5]$. A
parameter that reaches a bound has zero gradient and stays there; we report
which parameters do.

\subsection{Determinant part}
\label{ssec:determinant}

The single-determinant trial functions use $D[\phi_{1s}, \phi_{2s}]$, with
linear modulations, $c^{(2)}_{1s} = c^{(2)}_{2s} = 0$, and three variational
parameters, $\alpha_{2s}$, $c^{(1)}_{1s}$ and $c^{(1)}_{2s}$.

The multi-configuration trial functions use
\begin{equation}
\begin{aligned}
    D_{\mathrm{CAS}} = {} & D[\phi_{1s}, \phi_{2s}]
    + c_p \sum_{k} D[\phi_{1s}, \phi_{2p_k}] \\
    & + c_d \sum_{m} D[\phi_{1s}, \phi_{3d_m}] ,
\end{aligned}
    \label{eq:cas}
\end{equation}
with variational coefficients $c_p$ and $c_d$, both initialised at zero. The
three $2p_k$ terms, and the five $3d_m$ terms, enter with equal weights, so
that the state is a $^1S$ singlet. Two electrons in the $2s$ and $2p$
orbitals can form $^1S^e$ only as $2s^2$ or $2p^2$, since $2s2p$ has odd
parity. The first two terms of Eq.~\eqref{eq:cas} are therefore the complete
active space CAS(2,4) for this state. The third term adds the $3d^2$
configurations.

\subsection{Jastrow factor}
\label{ssec:jastrow}

When two electrons coalesce, the cusp condition reads~\cite{kato1957}
\begin{equation}
    \left. \frac{\partial J}{\partial r_{ij}} \right|_{r_{ij}\to 0} = a_{ij} ,
    \qquad
    a_{ij} =
    \begin{cases}
        1/4 & \text{parallel spins,}\\
        1/2 & \text{antiparallel spins.}
    \end{cases}
    \label{eq:ee_cusp}
\end{equation}
The two values differ because for parallel spins $D$ vanishes linearly at
coalescence. A Pad\'e term satisfies Eq.~\eqref{eq:ee_cusp} for any range
parameter,
\begin{equation}
    J_{\mathrm{P}} = \sum_{i<j} \frac{a_{ij}\, r_{ij}}{1 + b_{ij}\, r_{ij}} ,
    \label{eq:pade}
\end{equation}
with $b_{ij} = b_\parallel$ or $b_\perp$ according to the spins, bounded to
$[0.05, 10]$.

The multi-configuration trial functions add a three-body term,
\begin{equation}
    J = J_{\mathrm{P}} + \sum_{i<j} g_\theta(\mathbf{u}_{ij})\,
        e^{-0.3\,(t_i + t_j + t_{ij})} ,
    \label{eq:three_body}
\end{equation}
with $t_i = t(r_i)$, $t_{ij} = t(r_{ij})$, and the features
\begin{equation}
    \mathbf{u}_{ij} = \bigl( t_i + t_j,\; t_i t_j,\; t_{ij},\; s_{ij},\;
                      t_i t_j \cos\theta_{ij} \bigr) ,
    \label{eq:features}
\end{equation}
where $s_{ij} = \pm1$ for parallel and antiparallel spins and $\theta_{ij}$ is
the angle between $\mathbf{r}_i$ and $\mathbf{r}_j$. The features are
symmetric under $i \leftrightarrow j$. Since they are built on $t$, the term
does not change either cusp. $g_\theta$ is a multi-layer perceptron with two
hidden layers of 16 units and the SiLU activation, $\sigma(x) = x/(1+e^{-x})$,
whose second derivative does not vanish. The same weights are used for all
six pairs.

Some trial functions add a one-body term,
\begin{equation}
    J_1 = \sum_{i} w_\theta\bigl(\ln[1 + t(r_i)]\bigr) ,
    \label{eq:one_body}
\end{equation}
with $w_\theta$ a perceptron with one hidden layer of 8 SiLU units. $J_1$ has
zero slope at the nucleus, so Eq.~\eqref{eq:en_cusp} still holds. At large
$r$ it grows at most as $\ln r$, so $e^{J_1}$ is a power of $r$ and $\Psi$
stays normalisable.

\subsection{The trial functions}
\label{ssec:models}

\paragraph{Single-determinant series.}
The four trial functions share the determinant of Sec.~\ref{ssec:determinant}
and differ only in $J$. \mS{0} has $J = 0$. \mS{1} has the Pad\'e term,
Eq.~\eqref{eq:pade}. \mS{2} replaces it by a network
$f_\theta\bigl(r_{ij}/(1 + r_{ij}),\, s_{ij}\bigr)$, summed over pairs, which
can represent a cusp of any slope but must find the values of
Eq.~\eqref{eq:ee_cusp} by optimisation. \mS{3} adds to the Pad\'e term a
network acting on $t(r_{ij})/[1 + t(r_{ij})]$; since $t'(0) = 0$, this network
cannot change the slope at coalescence, and the cusps are those of the Pad\'e
term. The networks have two hidden layers of 16 units. The numbers of
variational parameters are 3, 5, 340 and 342.

\paragraph{Multi-configuration trial functions.}
Table~\ref{tab:models} lists the ten trial functions built on
Eq.~\eqref{eq:cas}. \mQ\ is the reference: all modulations are quadratic.
\mQb\ differs from \mQ\ only in the seed of the optimisation. \mL\ has
linear modulations, \mU\ none. \mQJ\ and \mUJ\ add the one-body term to \mQ\
and \mU. \mFJ\ keeps $m_{1s} = 1$, modulates $\phi_{2s}$ and $\phi_{2p}$, and
has the one-body term. \mD, \mDU\ and \mDJ\ add the $3d^2$ configurations to
\mQ, \mU\ and \mQJ. A parameter that is not used is kept at its initial value,
$c^{(1)} = c^{(2)} = 0$.

\begin{table}[t]
\caption{Multi-configuration trial functions. ``Modulation'' gives the form of
$m_{nl}$, Eq.~\eqref{eq:modulation}; $J_1$ is the one-body term,
Eq.~\eqref{eq:one_body}. The last column is the seed of the optimisation.}
\label{tab:models}
\centering
\small
\setlength{\tabcolsep}{3.5pt}
\begin{tabular}{lllccr}
\hline\hline
Label & Configurations & Modulation & $J_1$ & Param. & Seed \\
\hline
\mQ   & $2s^2, 2p^2$        & quadratic                & no  & 396 & 42 \\
\mQb  & $2s^2, 2p^2$        & quadratic                & no  & 396 & 7  \\
\mL   & $2s^2, 2p^2$        & linear                   & no  & 393 & 42 \\
\mU   & $2s^2, 2p^2$        & none                     & no  & 390 & 42 \\
\mQJ  & $2s^2, 2p^2$        & quadratic                & yes & 421 & 42 \\
\mUJ  & $2s^2, 2p^2$        & none                     & yes & 415 & 42 \\
\mFJ  & $2s^2, 2p^2$        & $2s$, $2p$; $m_{1s} = 1$ & yes & 419 & 42 \\
\mD   & $2s^2, 2p^2, 3d^2$  & quadratic                & no  & 400 & 42 \\
\mDU  & $2s^2, 2p^2, 3d^2$  & none                     & no  & 392 & 42 \\
\mDJ  & $2s^2, 2p^2, 3d^2$  & quadratic                & yes & 425 & 42 \\
\hline\hline
\end{tabular}
\end{table}


% ============================================================================
\section{Structure of the nodal surface}
\label{sec:nodes}
% ============================================================================

The Jastrow factor is positive, so the nodal surface of $\Psi$ is the nodal
surface of $D_{\mathrm{CAS}}$. In this section we write $D_{\mathrm{CAS}}$ as
the product of a positive function and a function that carries the node. We
then list the transformations of the orbitals that leave the node unchanged,
and the quantities that do not depend on them.

\subsection{Factorisation}
\label{ssec:factorisation}

For an orbital $\chi$, the spin-up determinant built from the pair
$(\phi_{1s}, \chi)$ can be written as
\begin{equation}
    \phi_{1s}(r_0)\,\phi_{1s}(r_1)\,
    \bigl[\tilde\chi(\mathbf{r}_1) - \tilde\chi(\mathbf{r}_0)\bigr] ,
    \qquad
    \tilde\chi = \frac{\chi}{\phi_{1s}} ,
    \label{eq:det_ratio}
\end{equation}
and the spin-down determinant in the same way. We define two radial functions,
\begin{align}
    f(r) &= \frac{\phi_{2s}(r)}{\phi_{1s}(r)}
          = (r - c)\, e^{(Z - \alpha_{2s}) r}\, q_{2s}(r) ,
    \label{eq:f} \\
    h(r) &= \frac{\phi_{2p_k}(\mathbf{r})}{x_k\, \phi_{1s}(r)}
          = e^{(Z - \alpha_{2p}) r}\, q_{2p}(r) ,
    \label{eq:h}
\end{align}
where $q_{nl} = m_{nl}/m_{1s}$. With $f_i = f(r_i)$ and
$\mathbf{w}_i = h(r_i)\,\mathbf{r}_i$, Eq.~\eqref{eq:cas} at $c_d = 0$ becomes
\begin{equation}
    D_{\mathrm{CAS}} = \Phi\, N ,
    \qquad
    \Phi = \prod_{i=0}^{3} \phi_{1s}(r_i) ,
    \label{eq:factorisation}
\end{equation}
\begin{equation}
\begin{aligned}
    N = {} & (f_1 - f_0)(f_3 - f_2) \\
      & + c_p\, (\mathbf{w}_1 - \mathbf{w}_0) \cdot (\mathbf{w}_3 - \mathbf{w}_2) .
\end{aligned}
    \label{eq:N}
\end{equation}
The sum over the three $2p_k^2$ configurations gives the scalar product.
$\Phi$ is positive as long as $m_{1s}$ has no zero; we checked this on a
radial grid up to $15$ bohr for every multi-configuration trial function. The
nodal surface is the set $N = 0$. The function $N$ depends on the electron
positions only through $f$, $h$ and $c_p$, id est through two radial
functions and one number.

The $3d^2$ configurations add one term of the same form. Summing over the
five configurations gives
\begin{equation}
\begin{aligned}
    N = {} & (f_1 - f_0)(f_3 - f_2) \\
      & + c_p\, (\mathbf{w}_1 - \mathbf{w}_0) \cdot (\mathbf{w}_3 - \mathbf{w}_2) \\
      & + c_d\, (W_1 - W_0) : (W_3 - W_2) ,
\end{aligned}
    \label{eq:N_d}
\end{equation}
with $W_i = h_d(r_i)\, T(\mathbf{r}_i)$,
$h_d(r) = e^{(Z - \alpha_{3d}) r}\, q_{3d}(r)$, and the colon denoting the full
contraction of the two tensors. A constant factor from the normalisation of
the $3d$ orbitals is absorbed in $c_d$.

\subsection{Redundancies}
\label{ssec:redundancies}

Two sets of transformations leave the node unchanged.

The first acts on $f$, $h$ and $c_p$:
\begin{equation}
    f \to a f + b , \qquad
    h \to \mu h , \qquad
    c_p \to c_p\, \frac{a^2}{\mu^2} ,
    \label{eq:gauge_psi}
\end{equation}
with constants $a \neq 0$, $b$ and $\mu \neq 0$. Under
Eq.~\eqref{eq:gauge_psi}, $N \to a^2 N$, so $\Psi$ changes by a constant
factor. The constant $b$ corresponds to adding a multiple of $\phi_{1s}$ to
$\phi_{2s}$. A determinant does not change when a multiple of one column is
added to another,
\begin{equation}
    \det\bigl(\phi_{1s},\, \phi_{2s} + \lambda\phi_{1s}\bigr)
    = \det\bigl(\phi_{1s},\, \phi_{2s}\bigr) ,
    \label{eq:column_invariance}
\end{equation}
so the orthogonalisation of $\phi_{2s}$ against $\phi_{1s}$ has no effect on
$\Psi$, and we omit it. The constants $a$ and $\mu$ are the scales of
$\phi_{2s}$ and $\phi_{2p}$ relative to $\phi_{1s}$. In the parametrisation of
Sec.~\ref{sec:wavefunction} these scales are not separate parameters, since
$m_{nl}(0) = 1$; the radial parameters change them, over the region occupied
by the electrons, together with the shape of the orbitals. The coefficient
$c_p$ therefore has no meaning on its own: two trial functions with nearly the
same node can have different $c_p$. The same holds for $c_d$, with
$W \to \nu W$ and $c_d \to c_d\, a^2/\nu^2$.

The second set acts on the modulations. If every $m_{nl}$ is multiplied by the
same positive function $g(r)$, $D_{\mathrm{CAS}}$ is multiplied by
$\prod_i g(r_i)$. This is a one-body Jastrow factor. The functions $f$ and $h$
do not change, and neither does the node. The node depends on the modulations
only through the ratios $q_{nl}$. Since $\alpha_{1s}$ is fixed at $Z$ by the
cusp condition, $m_{1s}$ is the only radial freedom of the $1s$ orbital, and it
enters the node only through these ratios.

\subsection{Invariants}
\label{ssec:invariants}

For a radius $r_\star$ we define
\begin{gather}
    \hat f(r) = \frac{f(r) - f(0)}{f(r_\star) - f(0)} , \qquad
    \hat h(r) = \frac{h(r)}{h(r_\star)} ,
    \nonumber \\
    \tilde c_p = c_p \left[ \frac{h(r_\star)}{f(r_\star) - f(0)} \right]^2 .
    \label{eq:invariants}
\end{gather}
These quantities do not change under Eq.~\eqref{eq:gauge_psi}, and
\begin{equation}
\begin{aligned}
    \frac{N}{\bigl[f(r_\star) - f(0)\bigr]^2} = {} &
        (\hat f_1 - \hat f_0)(\hat f_3 - \hat f_2) \\
        & + \tilde c_p\, (\hat{\mathbf{w}}_1 - \hat{\mathbf{w}}_0) \cdot
                       (\hat{\mathbf{w}}_3 - \hat{\mathbf{w}}_2) ,
\end{aligned}
    \label{eq:N_hat}
\end{equation}
with $\hat{\mathbf{w}}_i = \hat h(r_i)\, \mathbf{r}_i$. The triple
$(\hat f, \hat h, \tilde c_p)$ fixes the node. The number $\tilde c_p$ alone
does not: trial functions with different $\hat f$ or $\hat h$ have different
nodes at equal $\tilde c_p$, and the value of $\tilde c_p$ depends on
$r_\star$. A comparison of $\tilde c_p$ between trial functions is meaningful
only when their $\hat f$ and $\hat h$ are close over the region occupied by
the electrons.

To measure how close they are, we match each trial function to a reference
$R$. With $g = f - f(0)$ and a radial weight $\rho(r)$, the constants
\begin{gather}
    a_\star = \arg\min_a \int \! dr\, \rho\,
              \frac{(a g - g_R)^2}{g_R^2} ,
    \nonumber \\
    \mu_\star = \arg\min_\mu \int \! dr\, \rho\,
              \frac{(\mu h - h_R)^2}{h_R^2}
    \label{eq:match}
\end{gather}
bring $f$ and $h$ closest to those of $R$. The matched coefficient
$\tilde c_p^{R} = c_p\, a_\star^2/\mu_\star^2$ is invariant under
Eq.~\eqref{eq:gauge_psi} and has no free radius. The relative residuals of the
two fits, $\epsilon_f$ and $\epsilon_h$, measure how far the radial functions
of the trial function are from those of $R$.

We also use the weight of each configuration channel,
\begin{equation}
    w_l = \frac{\int \! d\mathbf{R}\; e^{2J}\, \Phi^2\, (c_l\, T_l)^2}
               {\int \! d\mathbf{R}\; e^{2J}\, \Phi^2\, \sum_{l'} (c_{l'}\, T_{l'})^2} ,
    \qquad l \in \{s, p, d\} ,
    \label{eq:weights}
\end{equation}
where $T_s$, $T_p$ and $T_d$ are the three products in Eq.~\eqref{eq:N_d} and
$c_s = 1$. The integrand of the denominator has no node, so $w_l$ is sampled
without difficulty. The weights do not change under
Eq.~\eqref{eq:gauge_psi}, since every channel scales by $a^2$. They depend on
the positive factor $e^{J}\Phi$, which the node does not fix: two trial
functions with the same node can have different $w_l$. The weights describe
the trial function, not the node.

\subsection{The limit $c_p = 0$}
\label{ssec:hf_limit}

At $c_p = c_d = 0$, $N = (f_1 - f_0)(f_3 - f_2)$. If $f$ is strictly monotonic
in $r$, $N$ vanishes exactly when $r_0 = r_1$ or $r_2 = r_3$. The node is then
\begin{equation}
    \{ r_0 = r_1 \} \cup \{ r_2 = r_3 \} ,
    \label{eq:hf_node}
\end{equation}
the node of the $1s^2 2s^2$ determinant~\cite{bressanini_nodes}, whatever the
radial parameters. We checked that $f$ is monotonic for every
multi-configuration trial function.
% TODO: verify monotonicity of f for the single-determinant series S0--S3.
The four trial functions of the single-determinant series therefore share one
node, and their energies differ only through the positive factor.

\subsection{Sign disagreement as a probe of the node}
\label{ssec:sign_probe}

A direct comparison of two nodal surfaces $a$ and $b$ is the probability that
the two trial functions have opposite signs,
\begin{equation}
    \varphi_{ab} = P_{|\Psi_a|^2}\bigl(\operatorname{sign}\Psi_a
                   \neq \operatorname{sign}\Psi_b\bigr) ,
    \label{eq:phi_ab}
\end{equation}
sampled from $|\Psi_a|^2$. This probe is weak. Near the node of $\Psi_a$,
$|\Psi_a|^2 \approx |\nabla\Psi_a|^2 d^2$, with $d$ the distance from the node.
If the node of $\Psi_b$ is displaced by $\delta(s)$ at the point $s$ of the
surface, the region between the two nodes carries the weight
\begin{equation}
    \varphi_{ab} \approx
    \frac{1}{3}\, \frac{\oint dS\; |\nabla\Psi_a|^2\, \delta^3}
                       {\int d\mathbf{R}\; |\Psi_a|^2} .
    \label{eq:phi_scaling}
\end{equation}
The probe is of third order in $\delta$. The fixed-node energy changes to
first order in $\delta$ for two generic nodes, and to second order about the
exact node. A sample without sign disagreements bounds $\varphi_{ab}$, but it
does not bound the difference of the fixed-node energies.


% ============================================================================
\section{Methods}
\label{sec:methods}
% ============================================================================

\subsection{Variational Monte Carlo}
\label{ssec:vmc_method}

We represent the trial function by $L = \ln|\Psi|$. The local energy is
\begin{equation}
    E_L = -\frac{1}{2}\bigl(\nabla^2 L + |\nabla L|^2\bigr)
          - \sum_i \frac{Z}{r_i} + \sum_{i<j} \frac{1}{r_{ij}} ,
    \label{eq:eloc}
\end{equation}
with the derivatives taken over the twelve coordinates by automatic
differentiation. Configurations are sampled from $|\Psi|^2$ with $4000$
Metropolis walkers. Each move displaces one electron by a Gaussian step; a
sweep moves the four electrons in turn. Single-electron moves are used
because the $1s$ and valence electrons have different length scales. The step
is adapted to an acceptance of about $55\%$.

The parameters are optimised by stochastic reconfiguration
(SR)~\cite{sorella1998}. With the scores
$O_k = \partial_{\theta_k} \ln|\Psi|$, the gradient of the energy and the
metric are
\begin{align}
    g_k &= 2\, \bigl\langle (O_k - \langle O_k\rangle)
                           (E_L - \langle E_L\rangle) \bigr\rangle ,
    \label{eq:gradient} \\
    S_{kl} &= \bigl\langle (O_k - \langle O_k\rangle)
                          (O_l - \langle O_l\rangle) \bigr\rangle .
    \label{eq:s_matrix}
\end{align}
The update is $\Delta\theta = -\eta\, (S + \gamma\,\mathrm{diag}\,S
+ 10^{-8})^{-1} g$, with $\gamma = 5\times10^{-3}$. The damping is relative
to the diagonal of $S$ because orbital parameters and network weights have
different scales. The norm of $\Delta\theta$ is then limited to
$\Delta_{\max}$. The learning rate and the trust radius start at
$\eta_0 = 0.1$ and $\Delta_0 = 0.3$ and decrease together as
$0.05 + 0.475\,[1 + \cos(\pi n/n_{\max})]$, with $n_{\max} = 3500$ epochs of
ten sweeps each. In the gradient only, the local energies are clipped to five
mean absolute deviations about their median. After the optimisation the
energy is accumulated over $200$ blocks separated by $150$ sweeps, and its
error is the largest standard error over the levels of a blocking
analysis~\cite{flyvbjerg1989}.

The single-determinant series was computed with an earlier version of the
code. It used the same sampler and SR update, with $\gamma$ between
$10^{-3}$ and $5\times10^{-3}$, $300$ (\mS{0}), $600$ (\mS{1}) and $1500$
(\mS{2}, \mS{3}) epochs, and $40$ blocks separated by $150$ sweeps.

\subsection{Fixed-node diffusion Monte Carlo}
\label{ssec:dmc_method}

FN-DMC samples the mixed distribution $\Psi\,\Phi_{\mathrm{FN}}$ with walkers
that drift, diffuse and carry weights~\cite{reynolds1982,foulkes2001}. A
walker moves as
$\mathbf{R}' = \mathbf{R} + \Delta\tau\, \bar{\mathbf{v}} + \bm{\eta}$, where
$\bm\eta$ is Gaussian with variance $\Delta\tau$ per coordinate and
$\bar{\mathbf{v}}$ is the drift $\nabla \ln|\Psi|$, reduced electron by
electron by the factor $2/(1 + \sqrt{1 + 2 v^2 \Delta\tau})$ of
Ref.~\onlinecite{umrigar1993}. The move is accepted with the Metropolis
probability built on the drift-diffusion Green function. A move that changes
the sign of $\Psi$ is rejected; this imposes the node. The weight is
multiplied by
$\exp\{-\Delta\tau\,[(E_L(\mathbf{R}) + E_L(\mathbf{R}'))/2 - E_T]\}$, with
$E_L$ limited to $\pm 2/\sqrt{\Delta\tau}$ about a reference energy. Every
five steps the weights are renormalised to a total of $N_w$, and the trial
energy is set to $E_T = E_{\mathrm{run}} - \ln(W/N_w)/\tau_{\mathrm{fb}}$
with $\tau_{\mathrm{fb}} = 1$. When the effective sample size falls below
$0.95\, N_w$ the population is resampled with a comb. The energy is the
weighted average of $E_L$.

All lengths are given in imaginary time. The walkers start from the VMC
walkers. After an equilibration of $\tau = 10$ the energy is accumulated in
blocks of $\tau = 1$ over $\tau = 150$, or $\tau = 100$ for the trial functions
with $J_1$. The time step is $\Delta\tau = 0.0025$ and $N_w = 4000$. Each
production run has its own seed, except \mQ\ and \mU, which share one; the
correlation of their blocks is $+0.04$, and we quote the error of their
difference without correction. The error is obtained by blocking with a
plateau test~\cite{flyvbjerg1989}. For every production run the plateau is
reached at the first level: blocks of $\tau = 1$ are uncorrelated.

\subsection{Nodal analysis}
\label{ssec:node_method}

The functions $f$, $h$, $\hat f$, $\hat h$ and the coefficients $\tilde c_p$
and $\tilde c_p^{R}$ follow from the parameters without sampling. We take
$r_\star = 2$ bohr unless stated otherwise. For Eq.~\eqref{eq:match} the
reference is \mQ\ and $\rho$ is the one-body radial density of \mQ, the same
for every trial function. Its 5\%, 50\% and 95\% quantiles are $0.14$, $0.95$
and $3.9$ bohr. The weights $w_l$ and the probabilities $\varphi_{ab}$ are
sampled with $4000$ Metropolis chains over $40$ blocks of $10$ sweeps; their
errors are the spread of the chain means divided by $\sqrt{4000}$.

\subsection{Tests}
\label{ssec:tests}

We tested the FN-DMC code on the hydrogen atom with the nodeless trial
function $e^{-1.3 r}$, at four time steps from $0.02$ to $0.0025$. The linear
extrapolation to $\Delta\tau = 0$ agrees with $-0.5\Eh$ within $0.4$ standard
deviations. This test has no node, and does not test the fixed-node
constraint. For \mQ\ we computed a time-step series, $\Delta\tau = 0.01$,
$0.005$ and $0.0025$ with $\tau = 40$ each, and runs with $1000$, $2000$ and
$4000$ walkers with $\tau = 60$ each.

The code is written in JAX~\cite{jax2018} with Flax, in double precision. A
production run of $\tau = 150$ with $4000$ walkers takes about $70$ minutes
on a laptop CPU.


% ============================================================================
\section{Results}
\label{sec:results}
% ============================================================================

\subsection{Single-determinant series}
\label{ssec:series}

Table~\ref{tab:series} gives the energies of the series. The four trial
functions have the same node (Sec.~\ref{ssec:hf_limit}), so the differences
come from the positive factor.
% TODO: this relies on f being monotonic for S0--S3 (not yet checked).

\begin{table*}[t]
\caption{Single-determinant series. Variational energy, variance of the local
energy, and fraction of the correlation energy, computed with
$E_{\mathrm{HF}} = -14.573023\Eh$~\cite{todohflimit} and
$E_{\mathrm{exact}} = -14.667356\Eh$~\cite{pachucki2006}. Bottom: slopes of
$J$ at coalescence compared with Eq.~\eqref{eq:ee_cusp}.}
\label{tab:series}
\centering
\small
\setlength{\tabcolsep}{4pt}
\begin{tabular}{lcccc}
\hline\hline
 & \mS{0} & \mS{1} & \mS{2} & \mS{3} \\
\hline
Parameters             & 3 & 5 & 340 & 342 \\
$E$ ($E_h$)            & $-14.5696(45)$ & $-14.6232(11)$
                       & $-14.6329(11)$ & $-14.6331(12)$ \\
$\mathrm{Var}(E_L)$    & 3.208 & 0.256 & 0.246 & 0.250 \\
Correlation            & $-3.6\%$ & $53.2\%$ & $63.5\%$ & $63.7\%$ \\
\hline
$\partial J/\partial r_{ij}|_0$, parallel ($1/4$)
                       & -- & 0.2500 & 0.4529 & 0.2500 \\
$\partial J/\partial r_{ij}|_0$, antiparallel ($1/2$)
                       & -- & 0.5000 & 0.4964 & 0.5000 \\
\hline\hline
\end{tabular}
\end{table*}

\mS{0} lies $3.4\mHa$ above the Hartree-Fock limit. Its determinant has three
radial parameters, and the Hartree-Fock limit is a lower bound for it. The
Pad\'e term, \mS{1}, lowers the energy by $53.6\mHa$ and the variance by a
factor of $12.5$. The network of \mS{2} lowers the energy by a further
$9.7\mHa$ and the variance by $4\%$. Adding the exact cusps to the network,
\mS{3}, changes the energy by $0.2(1.6)\mHa$.

The network of \mS{2} reproduces the antiparallel-spin cusp to $0.7\%$ and
gives the parallel-spin pairs a slope of $0.4529$ instead of $0.25$, an error
of $81\%$. For parallel spins $D$ vanishes at coalescence, so the region
$r_{ij} \to 0$ carries little weight for these pairs, and only two of the six
pairs have parallel spins. The gradient of the energy along the parallel-spin
slope is therefore small. \mS{3} returns the exact slopes to four decimals,
as required by its construction.

Figure~\ref{fig:jastrow} shows the optimised Jastrow factors. In \mS{1} the
parallel-spin range parameter runs to the bound, $b_\parallel = 10$: the
parallel-spin term becomes a constant, $a/b$, and is switched off. With the
bound raised to $100$, $b_\parallel$ runs again to the bound and the energy
rises to $-14.5787\Eh$. In \mS{3}, $b_\parallel = 0.283$, inside the interval.
\mS{1} and \mS{3} have opposite parallel-spin correlation functions and the
same energy within the error. In \mS{2} the network grows linearly at large
$r_{ij}$, with slopes $0.088$ (antiparallel) and $0.07$ (parallel). For an
electron far from the others this gives a factor of about $e^{0.25 r}$, which
compensates part of the orbital decay; $\alpha_{2s}$ is $1.314$ in \mS{2}
against $1.170$ in \mS{1}. The network does one-body work, and the network and
$\alpha_{2s}$ move along a direction of nearly constant energy.

\begin{figure*}[t]
\centering
\figfile{\textwidth}{comparison_jastrow.pdf}
\caption{Optimised Jastrow exponents of \mS{1}, \mS{2} and \mS{3}, shifted to
$J(0) = 0$. Solid lines: antiparallel pairs; dashed lines: parallel pairs.
Shaded histogram: distribution of the sampled $r_{ij}$ (right axis). Right
panel: the coalescence region, with the exact cusp slopes $1/4$ and $1/2$ as
dotted lines.}
\label{fig:jastrow}
\end{figure*}

\subsection{Multi-configuration trial functions: VMC}
\label{ssec:vmc}

Table~\ref{tab:vmc} gives the VMC results. The energies lie between
$-14.66620\Eh$ (\mU) and $-14.66697\Eh$ (\mDJ). The errors are
$0.11$--$0.17\mHa$, and about half of the differences between trial
functions are below $1.5$ standard deviations. The variance of the local energy separates them
more: $0.0086$ for \mQJ\ and \mDJ, $0.0095$ for \mQ, $0.0131$ for \mU.

\begin{table}[t]
\caption{VMC energies of the multi-configuration trial functions, variance of
the local energy, fraction of the correlation energy, and mixing coefficient
$c_p$.}
\label{tab:vmc}
\centering
\small
\setlength{\tabcolsep}{4pt}
\begin{tabular}{lcccc}
\hline\hline
 & $E_{\mathrm{VMC}}$ ($E_h$) & $\mathrm{Var}(E_L)$ & Corr. (\%) & $c_p$ \\
\hline
\mQ   & $-14.66677(12)$ & 0.0095 & 99.38(13) & $-0.2200$ \\
\mQb  & $-14.66669(17)$ & 0.0109 & 99.29(18) & $-0.2226$ \\
\mL   & $-14.66660(12)$ & 0.0113 & 99.19(13) & $-0.1984$ \\
\mU   & $-14.66620(13)$ & 0.0131 & 98.77(14) & $-0.2437$ \\
\mQJ  & $-14.66689(11)$ & 0.0086 & 99.51(11) & $-0.2184$ \\
\mUJ  & $-14.66633(14)$ & 0.0122 & 98.91(15) & $-0.2372$ \\
\mFJ  & $-14.66676(14)$ & 0.0115 & 99.36(14) & $-0.1904$ \\
\mD   & $-14.66663(12)$ & 0.0107 & 99.23(13) & $-0.2387$ \\
\mDU  & $-14.66647(13)$ & 0.0134 & 99.06(14) & $-0.2569$ \\
\mDJ  & $-14.66697(12)$ & 0.0086 & 99.59(13) & $-0.2376$ \\
\hline\hline
\end{tabular}
\end{table}

Some coefficients end at a bound: $c^{(2)}_{1s} = 0.8$ in \mQJ\ and \mDJ, both
coefficients of $m_{3d}$ in \mD\ and \mDJ, and $c^{(1)}_{2s} = -0.15$ in \mFJ.
For these trial functions the results depend on the bounds. In every trial
function $m_{1s}$ has no zero and $f$ is monotonic. Equation~\eqref{eq:N_d}
reproduces the directly computed determinant part to a relative accuracy of
$10^{-14}$.

In \mQJ\ the one-body term is negative: $w_\theta(r) - w_\theta(0) = -0.55$,
$-1.10$ and $-1.79$ at $r = 1$, $2$ and $4$ bohr. In the same trial function
$m_{1s}$ grows more than in \mQ: $c^{(1)}_{1s} = 1.55$ against $1.30$, and
$c^{(2)}_{1s}$ reaches its bound. At $r = 4$ bohr, $\ln m_{1s} = 2.65$. The
one-body term removes part of the one-body factor that $m_{1s}$ carries. In
\mUJ, where the orbitals are not modulated, the one-body term is positive
($+0.51$ at $r = 4$ bohr) and $\alpha_{2p}$ increases from $1.10$ to $1.21$.

\subsection{Error budget}
\label{ssec:budget}

Table~\ref{tab:dmc} gives the FN-DMC energies. Figure~\ref{fig:budget}
compares $E_{\mathrm{VMC}}$ and $E_{\mathrm{FN}}$ for \mQ, \mL\ and \mU. In
every trial function $E_{\mathrm{VMC}} - E_{\mathrm{FN}}$ is larger than
$E_{\mathrm{FN}} - E_{\mathrm{exact}}$. For \mQ, \mQb, \mQJ\ and \mDJ\ the
positive factor accounts for $60$--$81\%$ of the VMC error, for \mFJ\ for
$52\%$ and for \mD\ for $98\%$. These fractions have errors of about $20$
percentage points.

\begin{table}[t]
\caption{FN-DMC energies at $\Delta\tau = 0.0025$, distance from the exact
energy, difference from the VMC energy, and fraction of the VMC error due to
the positive factor, $(E_{\mathrm{VMC}} - E_{\mathrm{FN}})/
(E_{\mathrm{VMC}} - E_{\mathrm{exact}})$. Energy differences in mHa.}
\label{tab:dmc}
\centering
\small
\setlength{\tabcolsep}{2.5pt}
\begin{tabular}{lcccc}
\hline\hline
 & $E_{\mathrm{FN}}$ ($E_h$) & $E_{\mathrm{FN}} - E_{\mathrm{ex}}$
 & $E_{\mathrm{VMC}} - E_{\mathrm{FN}}$ & Fraction \\
\hline
\mQ   & $-14.667237(37)$ & $+0.119(37)$ & $0.46(13)$ & 80\% \\
\mQb  & $-14.667204(33)$ & $+0.152(33)$ & $0.52(17)$ & 77\% \\
\mL   & $-14.667169(35)$ & $+0.187(35)$ & $0.57(13)$ & 75\% \\
\mU   & $-14.666973(43)$ & $+0.383(43)$ & $0.78(14)$ & 67\% \\
\mQJ  & $-14.667269(41)$ & $+0.087(41)$ & $0.38(11)$ & 81\% \\
\mUJ  & $-14.667045(58)$ & $+0.311(58)$ & $0.71(15)$ & 70\% \\
\mFJ  & $-14.667066(49)$ & $+0.290(49)$ & $0.31(15)$ & 52\% \\
\mD   & $-14.667340(69)$ & $+0.016(69)$ & $0.71(14)$ & 98\% \\
\mDJ  & $-14.667201(35)$ & $+0.155(35)$ & $0.23(13)$ & 60\% \\
\hline\hline
\end{tabular}
\end{table}

\begin{figure}[t]
\centering
\figfile{\columnwidth}{fig_error_budget_quadA.pdf}\\[2pt]
\figfile{\columnwidth}{fig_error_budget_lin.pdf}\\[2pt]
\figfile{\columnwidth}{fig_error_budget_nomod.pdf}
\caption{VMC and FN-DMC energies relative to the exact energy for \mQ\ (top),
\mL\ (middle) and \mU\ (bottom). The arrow is
$E_{\mathrm{VMC}} - E_{\mathrm{FN}}$, the error of the positive factor.}
\label{fig:budget}
\end{figure}

\subsection{Radial modulation and the node}
\label{ssec:modulation}

The fixed-node energy of \mQ\ is lower than that of \mU\ by $0.264(57)\mHa$,
$4.7$ standard deviations. The radial modulation changes the node. In VMC the
difference is $0.577(179)\mHa$. Of this, $0.264\mHa$ comes from the node and
$0.313(188)\mHa$ from the positive factor. The VMC energy alone does not
separate the two parts.

\mL\ lies between the two: $E_{\mathrm{FN}}(\mL) - E_{\mathrm{FN}}(\mU) =
-0.196(55)\mHa$ and $E_{\mathrm{FN}}(\mL) - E_{\mathrm{FN}}(\mQ) =
+0.068(51)\mHa$. \mQb, which differs from \mQ\ only in the seed of the
optimisation, gives $E_{\mathrm{FN}}(\mQb) - E_{\mathrm{FN}}(\mQ) =
+0.033(50)\mHa$. A second pair that differs in an element without effect on
the node, \mUJ\ and \mU\ (Sec.~\ref{ssec:dchannel}), gives $-0.072(72)\mHa$.
In both cases the spread is consistent with the statistical error.

\subsection{The node and the $1s$ orbital}
\label{ssec:onebody}

Table~\ref{tab:groups} divides the nine trial functions with FN-DMC energies
by construction: $m_{1s}$ optimised or $m_{1s} = 1$. Within each group the
fixed-node energies agree within their errors. The trial functions with
$m_{1s} = 1$ lie $0.198(32)\mHa$ higher. Figure~\ref{fig:groups} shows the
nine energies.

\begin{table}[t]
\caption{Fixed-node energies grouped by the treatment of the $1s$ orbital.
Weighted means, in mHa above the exact energy, and $\chi^2$ about the mean.}
\label{tab:groups}
\centering
\small
\setlength{\tabcolsep}{4pt}
\begin{tabular}{llcc}
\hline\hline
Group & Trial functions & $E_{\mathrm{FN}} - E_{\mathrm{ex}}$ & $\chi^2$/dof \\
\hline
$m_{1s}$ optimised & \parbox[t]{2.6cm}{\raggedright \mQ, \mQb, \mL, \mQJ, \mD, \mDJ} & $+0.137(16)$ & 7.3/5 \\
$m_{1s} = 1$       & \mU, \mUJ, \mFJ               & $+0.335(28)$ & 2.3/2 \\
Difference         &                                & $0.198(32)$  &       \\
\hline\hline
\end{tabular}
\end{table}

\begin{figure}[t]
\centering
\figfile{\columnwidth}{fig_node_groups.pdf}
% TODO: figure to be produced. E_FN - E_exact of the nine trial functions,
% grouped by m_1s optimised / m_1s = 1, with the two weighted means as bands.
\caption{Fixed-node energies of the nine trial functions relative to the exact
energy, grouped by the treatment of the $1s$ orbital. Bands: weighted means of
Table~\ref{tab:groups}.}
\label{fig:groups}
\end{figure}

The groups also differ in other elements: \mU\ and \mUJ\ have no modulation
of $\phi_{2s}$ and $\phi_{2p}$. Three trial functions with the same Jastrow
factor isolate the $1s$ orbital. From \mUJ\ to \mFJ, $\phi_{2s}$ and
$\phi_{2p}$ are modulated and $\phi_{1s}$ is not; the fixed-node energy
changes by $-0.02(8)\mHa$. From \mFJ\ to \mQJ, $m_{1s}$ is added; the
fixed-node energy changes by $-0.20(6)\mHa$, $3.2$ standard deviations. In
VMC, \mQJ\ and \mFJ\ differ by $0.14(17)\mHa$. Prima facie the two trial
functions are equivalent.

The modulation of the $1s$ orbital acts on the node through the ratios
$q_{nl}$ of Sec.~\ref{ssec:redundancies}. With $\alpha_{1s}$ fixed at $Z = 4$
by the cusp condition, $m_{1s}$ is the only way to change the shape of
$\phi_{1s}$. In \mQ\ and \mQJ\ it grows with $r$ and makes $\phi_{1s}$ more
diffuse than $e^{-Zr}$. Since $f = \phi_{2s}/\phi_{1s}$ and
$h = \phi_{2p_k}/(x_k\phi_{1s})$, the shape of $\phi_{1s}$ enters the node.
In \mFJ\ the modulations of $\phi_{2s}$ and $\phi_{2p}$ cannot replace it:
$c^{(1)}_{2s}$ ends at its lower bound. In \mQJ\ the coefficient
$c^{(2)}_{1s}$ is at its upper bound, so the size of the effect may depend on
the bound.

\subsection{Redundancies in the optimised parameters}
\label{ssec:redundancies_results}

Table~\ref{tab:nodal} gives, for every trial function, $c_p$, $\tilde c_p$ at
$r_\star = 2$ bohr, the matched coefficient $\tilde c_p^{R}$ with \mQ\ as
reference, and the residuals of the match. Figure~\ref{fig:radial} shows the
modulations, the ratios $q_{nl}$ and the functions $\hat f$, $\hat h$ for \mQ,
\mQb, \mL\ and \mU.

\begin{table}[t]
\caption{Mixing coefficient $c_p$, invariant $\tilde c_p$ at
$r_\star = 2$ bohr, matched coefficient $\tilde c_p^{R}$
(Eq.~\eqref{eq:match}, reference \mQ), and relative residuals of the match.}
\label{tab:nodal}
\centering
\small
\setlength{\tabcolsep}{3.5pt}
\begin{tabular}{lccccc}
\hline\hline
 & $c_p$ & $\tilde c_p$ & $\tilde c_p^{R}$ & $\epsilon_f$ & $\epsilon_h$ \\
\hline
\mQ   & $-0.2200$ & $-0.1078$ & $-0.2200$ & 0     & 0     \\
\mQb  & $-0.2226$ & $-0.1084$ & $-0.2211$ & 0.024 & 0.024 \\
\mL   & $-0.1984$ & $-0.1060$ & $-0.2115$ & 0.156 & 0.168 \\
\mU   & $-0.2437$ & $-0.0998$ & $-0.1936$ & 0.639 & 0.607 \\
\mQJ  & $-0.2184$ & $-0.1081$ & $-0.2121$ & 0.043 & 0.030 \\
\mUJ  & $-0.2372$ & $-0.1001$ & $-0.1947$ & 0.547 & 0.491 \\
\mFJ  & $-0.1904$ & $-0.1058$ & $-0.2174$ & 0.483 & 0.460 \\
\mD   & $-0.2387$ & $-0.1091$ & $-0.2301$ & 0.062 & 0.071 \\
\mDU  & $-0.2569$ & $-0.1015$ & $-0.1888$ & 0.629 & 0.586 \\
\mDJ  & $-0.2376$ & $-0.1080$ & $-0.2212$ & 0.063 & 0.071 \\
\hline\hline
\end{tabular}
\end{table}

\begin{figure*}[t]
\centering
\figfile{\textwidth}{fig_gauge_radial.pdf}
\caption{Radial functions of \mQ, \mQb, \mL\ and \mU. (a)--(c) Modulations
$m_{1s}$, $m_{2s}$, $m_{2p}$. (d), (e) Ratios $q_{2s}$ and $q_{2p}$, which
enter the node. (f) $\hat f$ and $\hat h$ at $r_\star = 2$ bohr, as ratios to
those of \mQ; by construction the ratios equal one at $r_\star$. Shaded band:
central 90\% of the sampled electron radii.}
\label{fig:radial}
\end{figure*}

Among the nine trial functions with FN-DMC energies the raw coefficient $c_p$
ranges from $-0.190$ to $-0.244$. It does not order the fixed-node energies: a linear fit of $E_{\mathrm{FN}}$ in $|c_p|$ gives
$\chi^2 = 46.8$ for $7$ degrees of freedom. The five trial functions whose
radial functions lie within $7.1\%$ of those of \mQ\ (\mQ, \mQb, \mQJ, \mD,
\mDJ) have $c_p$ between $-0.218$ and $-0.239$ and fixed-node energies that
agree within their errors ($\chi^2 = 4.8/4$). For these trial functions the
spread of $\tilde c_p^{R}$ is $8\%$.

The trial functions with $m_{1s} = 1$ have radial functions $46$--$64\%$ away
from those of \mQ. For them a single number cannot describe the node.
Table~\ref{tab:rstar} shows the consequence: a linear fit of
$E_{\mathrm{FN}}$ in $|\tilde c_p|$ over the nine trial functions has a slope
that changes sign with $r_\star$. At small $r_\star$ the trial functions with
$m_{1s} = 1$ have the largest $|\tilde c_p|$, at large $r_\star$ the smallest.
The fit separates the two groups of Table~\ref{tab:groups} in either
direction. The matched coefficient gives $\chi^2 = 17.8/7$, and no slope
within the five modulated trial functions without $3d^2$ configurations.

\begin{table}[t]
\caption{Linear fit of $E_{\mathrm{FN}} - E_{\mathrm{ex}}$ in
$|\tilde c_p|$ for the nine trial functions, at several gauge-fixing radii and
with the matched coefficient. Slope in units of its standard error.}
\label{tab:rstar}
\centering
\small
\begin{tabular}{lcc}
\hline\hline
Gauge & Slope ($\sigma$) & $\chi^2$/dof \\
\hline
$r_\star = 0.5$ bohr & $+5.7$ & 14.1/7 \\
$r_\star = 1.0$ bohr & $+5.3$ & 18.9/7 \\
$r_\star = 1.5$ bohr & $-3.0$ & 37.9/7 \\
$r_\star = 2.0$ bohr & $-6.2$ & 8.7/7  \\
$r_\star = 3.0$ bohr & $-5.6$ & 16.4/7 \\
Matched to \mQ       & $-5.4$ & 17.8/7 \\
\hline\hline
\end{tabular}
\end{table}

The coefficient of the $3d^2$ configurations behaves in the same way. From
\mDU\ to \mD, $|c_d|$ decreases from $0.078$ to $0.016$, while the weight of
the channel, Eq.~\eqref{eq:weights}, increases from $6.8\times10^{-5}$ to
$1.05\times10^{-4}$. The change of $c_d$ comes from the scale of the $3d$
orbital: in \mD, $\alpha_{3d} = 2.02$ and both coefficients of $m_{3d}$ are at
their bounds.

The VMC energy does not order the fixed-node energies either. The rank
correlation of $E_{\mathrm{VMC}}$ and $E_{\mathrm{FN}}$ over the nine trial
functions is $0.57$ ($p = 0.11$). For the variance of the local energy it is
$0.81$ ($p = 0.008$).

\subsection{The $3d^2$ configurations and the one-body term}
\label{ssec:dchannel}

Two pairs of trial functions differ only in the $3d^2$ configurations:
$E_{\mathrm{FN}}(\mD) - E_{\mathrm{FN}}(\mQ) = -0.103(78)\mHa$ and
$E_{\mathrm{FN}}(\mDJ) - E_{\mathrm{FN}}(\mQJ) = +0.068(54)\mHa$. The weighted
mean is $+0.01(4)\mHa$. The $3d^2$ configurations do not change the
fixed-node energy within $0.04\mHa$, in agreement with a channel weight of
$10^{-4}$. \mD\ lies $1.8$ standard deviations below the mean of the other
five trial functions with optimised $m_{1s}$. Its error is twice that of the
other runs of the same length. The standard deviation of its block energies is $0.84\mHa$, against
$0.35$--$0.46\mHa$ for \mQ, \mQb\ and \mDJ. Its excess kurtosis, $1.5$, lies
within the range of the other runs ($-0.1$ to $2.8$), so the width does not
come from a few outlying blocks. The blocks are uncorrelated, and the error is
the standard error of the block mean.

Three pairs differ only in the one-body term:
$E_{\mathrm{FN}}(\mQJ) - E_{\mathrm{FN}}(\mQ) = -0.032(56)\mHa$,
$E_{\mathrm{FN}}(\mUJ) - E_{\mathrm{FN}}(\mU) = -0.072(72)\mHa$ and
$E_{\mathrm{FN}}(\mDJ) - E_{\mathrm{FN}}(\mD) = +0.139(77)\mHa$. The weighted
mean is $0.00(4)\mHa$. In VMC the one-body term lowers the energy of \mU\ by
$0.14(19)\mHa$ and that of \mQ\ by $0.12(16)\mHa$. It lowers the variance of
\mQ\ and \mU\ by $7$--$10\%$ and that of \mD\ by $20\%$.

\subsection{Sign disagreement}
\label{ssec:sign_results}

Among \mQ, \mQb, \mL\ and \mU, no pair has a sign disagreement in the sampled
configurations. With $4000$ independent chains this gives
$\varphi_{ab} < 7.5\times10^{-4}$ at $95\%$ confidence. The fixed-node
energies of \mQ\ and \mU\ differ by $0.264(57)\mHa$. Both facts follow from
Eq.~\eqref{eq:phi_scaling}: the probe is of third order in the displacement of
the node, the energy of first order.

\subsection{Systematic errors}
\label{ssec:systematics}

Table~\ref{tab:timestep} and Fig.~\ref{fig:timestep} give the time-step series
of \mQ. The data do not fix the form of the time-step error. A linear and a
quadratic fit extrapolate to $+0.19$ and $+0.15\mHa$ above the exact energy;
a constant fits the data with $\chi^2 = 6.8/3$. The production value at
$\Delta\tau = 0.0025$ is $+0.119(37)\mHa$. A linear or quadratic form places
the error at $\Delta\tau = 0.0025$ between $-0.05$ and $-0.01\mHa$. We quote
all fixed-node energies at $\Delta\tau = 0.0025$, with this range as a
systematic uncertainty. The series exists only for \mQ; for the differences
between trial functions we assume that the time-step error at
$\Delta\tau = 0.0025$ is the same within this range. The fraction of rejected
node crossings is $0.0169\%$, $0.0059\%$ and $0.0020\%$ at the three time
steps, and $0.0020$--$0.0022\%$ in every production run.

\begin{table}[t]
\caption{Time-step series of \mQ\ (top) and fits to it (bottom).
$E_{\mathrm{FN}} - E_{\mathrm{ex}}$ in mHa.}
\label{tab:timestep}
\centering
\small
\begin{tabular}{lcc}
\hline\hline
$\Delta\tau$ & $E_{\mathrm{FN}} - E_{\mathrm{ex}}$ & Crossings \\
\hline
$0.01$             & $-0.046(64)$ & $0.0169\%$ \\
$0.005$            & $+0.154(63)$ & $0.0059\%$ \\
$0.0025$           & $+0.152(65)$ & $0.0020\%$ \\
$0.0025$, production & $+0.119(37)$ & $0.0020\%$ \\
\hline
Fit & Intercept & $\chi^2$/dof \\
\hline
Constant  & $+0.103(26)$ & 6.8/3 \\
Linear    & $+0.193(48)$ & 1.7/2 \\
Quadratic & $+0.150(33)$ & 1.0/2 \\
\hline\hline
\end{tabular}
\end{table}

\begin{figure}[t]
\centering
\figfile{\columnwidth}{fig_dmc_timestep_quadA.pdf}
% TODO: the current figure draws only the constant fit; add the linear and
% quadratic fits before submission.
\caption{Time-step series of \mQ. (a) $E_{\mathrm{FN}} - E_{\mathrm{exact}}$
and fits. (b) Fraction of rejected node crossings, logarithmic scales.}
\label{fig:timestep}
\end{figure}

With $1000$, $2000$ and $4000$ walkers the energies of \mQ\ are
$-14.667378(107)$, $-14.667182(81)$ and $-14.667201(53)\Eh$. A constant fits
them with $\chi^2 = 2.5/2$; a term in $1/N_w$ has a slope of
$-0.21(16)\mHa$ at $N_w = 1000$. The population error at $N_w = 4000$ is
below the statistical error.

\subsection{Comparison with previous work}
\label{ssec:literature}

Table~\ref{tab:literature} compares our results with those of Seth
\emph{et al.}~\cite{seth2011} and Morales \emph{et al.}~\cite{morales2012}.
The VMC energy of \mDJ\ lies $0.22(12)\mHa$ above that of Seth \emph{et al.},
who use a larger expansion and backflow, and agrees with that of Morales
\emph{et al.}. The variance of \mDJ, $0.0086$, lies between the values of Seth
\emph{et al.} without backflow, $0.0107$, and with backflow, $0.0053$. The
fixed-node energy of the trial functions with optimised $m_{1s}$ lies
$0.087(17)\mHa$ above that of Seth \emph{et al.} Our value is not
extrapolated in $\Delta\tau$. The extrapolation would raise it by up to
$0.05\mHa$.
% TODO: verify every literature value against the papers.

\begin{table}[t]
\caption{Comparison with previous work. Energies in $E_h$.}
\label{tab:literature}
\centering
\small
\setlength{\tabcolsep}{2pt}
\begin{tabular}{lcc}
\hline\hline
 & $E_{\mathrm{VMC}}$ & $E_{\mathrm{FN}}$ \\
\hline
This work, \mDJ              & $-14.66697(12)$ & $-14.667201(35)$ \\
This work, mean, $m_{1s}$ opt. & --             & $-14.667219(16)$ \\
Seth \emph{et al.}~\cite{seth2011}     & $-14.66719(1)$ & $-14.667306(7)$ \\
Morales \emph{et al.}~\cite{morales2012} & $-14.66688(4)$ & -- \\
Exact~\cite{pachucki2006}    & \multicolumn{2}{c}{$-14.667356$} \\
\hline\hline
\end{tabular}
\end{table}


% ============================================================================
\section{Discussion}
\label{sec:discussion}
% ============================================================================

The network learns the parts of the Jastrow factor to which the energy is
sensitive. In the single-determinant series the network reproduces the antiparallel-spin cusp
and not the parallel-spin cusp, because the energy depends on the first and
hardly on the second. Imposing both cusps costs one change of variable and
gives the same energy. In the multi-configuration trial functions the positive
factor still carries most of the VMC error: $E_{\mathrm{VMC}} -
E_{\mathrm{FN}}$ is $0.2$--$0.8\mHa$ against a fixed-node error of
$0.02$--$0.38\mHa$.

The Jastrow factor cannot learn the node, since $e^J > 0$. The node is fixed
by the determinant part, and in the trial functions studied here by two
radial functions and one number. The parameters of the determinant part,
optimised together with the Jastrow factor, are not identifiable: $c_p$ and
$c_d$ change with the scale of the orbitals, and a common radial factor of the
orbitals is a one-body Jastrow factor. What changes the node is the shape of
each orbital relative to the $1s$ orbital. In our parametrisation the cusp
condition fixes the $1s$ exponent, and the modulation $m_{1s}$ is the radial
freedom that the node needs.

Neither the VMC energy nor the raw coefficients detect a change of the node.
\mFJ\ and \mQJ\ have the same VMC energy within $0.14(17)\mHa$ and fixed-node
energies $0.20(6)\mHa$ apart. A gauge-invariant coefficient detects the change
only between trial functions with the same radial functions; across different
radial functions its ordering depends on the choice of gauge. To compare nodes
of different radial form we used FN-DMC. The variance of the local energy
correlates with the fixed-node energy better than the VMC energy, but it also
depends on the positive factor.


% ============================================================================
\section{Limitations and outlook}
\label{sec:limitations}
% ============================================================================

The study treats one atom with a quasi-degenerate ground state. The effects on
the node are $0.1$--$0.3\mHa$, a few standard deviations of our fixed-node
energies. Only one pair of trial functions, \mQ\ and \mQb, differs only in the
seed of the optimisation. The time-step series and the population test exist
only for \mQ. The FN-DMC code was tested only on a nodeless system. Some
coefficients of the best trial functions end at their bounds. The
single-determinant series was computed with an earlier version of the code and
a shorter protocol.

The invariants of Sec.~\ref{sec:nodes} hold for the CAS(2,4) expansion with
the orbitals of Sec.~\ref{sec:wavefunction}. A measure of the distance between
two nodes that does not depend on a choice of gauge would allow nodes of
different radial form to be compared without FN-DMC. A parametrisation of the
ratios $q_{nl}$ instead of the modulations $m_{nl}$ would remove the
redundancy of Sec.~\ref{ssec:redundancies} and the bounds that now limit the
$1s$ orbital. The same analysis applies to other atoms with a
$2s^2$--$2p^2$ near-degeneracy.


% ============================================================================
% BIBLIOGRAPHY: placeholders, to be completed and checked.
% ============================================================================
\begin{thebibliography}{99}

\bibitem{carleo2017} TODO: G. Carleo and M. Troyer, Science (2017).
\bibitem{pfau2020} TODO: D. Pfau \emph{et al.}, Phys. Rev. Res. (2020), FermiNet.
\bibitem{hermann2020} TODO: J. Hermann, Z. Sch\"atzle and F. No\'e, Nat. Chem. (2020), PauliNet.
\bibitem{reynolds1982} TODO: P. J. Reynolds \emph{et al.}, J. Chem. Phys. (1982).
\bibitem{foulkes2001} TODO: W. M. C. Foulkes \emph{et al.}, Rev. Mod. Phys. (2001).
\bibitem{todobehfnode} TODO: source for the fixed-node energy of Be with a single determinant.
\bibitem{pachucki2006} TODO: K. Pachucki and J. Komasa, exact non-relativistic energy of Be.
\bibitem{bressanini_nodes} TODO: D. Bressanini, D. M. Ceperley and P. J. Reynolds, nodes of Be.
\bibitem{kato1957} TODO: T. Kato, Commun. Pure Appl. Math. (1957).
\bibitem{sorella1998} TODO: S. Sorella, Phys. Rev. Lett. (1998).
\bibitem{flyvbjerg1989} TODO: H. Flyvbjerg and H. G. Petersen, J. Chem. Phys. (1989).
\bibitem{umrigar1993} TODO: C. J. Umrigar, M. P. Nightingale and K. J. Runge, J. Chem. Phys. (1993).
\bibitem{jax2018} TODO: J. Bradbury \emph{et al.}, JAX (2018).
\bibitem{todohflimit} TODO: source for the Hartree-Fock limit of Be.
\bibitem{seth2011} TODO: P. Seth, P. L\'opez R\'ios and R. J. Needs, J. Chem. Phys. (2011).
\bibitem{morales2012} TODO: M. A. Morales \emph{et al.}, J. Chem. Theory Comput. (2012).

\end{thebibliography}

\end{document}